% ============================================================
%  Глава 0 — Теоретические основы теории управления
%  Справочная глава: модели, линеаризация, устойчивость
% ============================================================
\chapter{Теоретические основы систем управления}
\label{ch:theory}

Данная глава содержит математический аппарат теории автоматического
управления (ТАУ), на котором базируются все алгоритмы проекта SAMURAI.

% ═══════════════════════════════════════════════════════════════
\section{Модель в пространстве состояний}
\label{sec:state_space}

Любой объект управления может быть описан системой дифференциальных
уравнений первого порядка --- \textbf{уравнениями состояния}:

\begin{equation}
  \boxed{
  \dot{\vect{x}}(t) = \mat{A}\,\vect{x}(t) + \mat{B}\,\vect{u}(t), \qquad
  \vect{y}(t) = \mat{C}\,\vect{x}(t) + \mat{D}\,\vect{u}(t)
  }
  \label{eq:state_space}
\end{equation}

где:
\begin{itemize}
  \item $\vect{x}(t) = (x_1, x_2, \ldots, x_n)\T \in \mathbb{R}^n$ --- \textbf{вектор состояния}
  \item $\vect{u}(t) = (u_1, \ldots, u_r)\T \in \mathbb{R}^r$ --- вектор управляющих воздействий
  \item $\vect{y}(t) = (y_1, \ldots, y_m)\T \in \mathbb{R}^m$ --- вектор выходных переменных
  \item $\mat{A} \in \mathbb{R}^{n \times n}$ --- матрица системы (динамики)
  \item $\mat{B} \in \mathbb{R}^{n \times r}$ --- матрица управления
  \item $\mat{C} \in \mathbb{R}^{m \times n}$ --- матрица наблюдения
  \item $\mat{D} \in \mathbb{R}^{m \times r}$ --- матрица прямого прохода
\end{itemize}

\begin{notebox}
Структурная схема модели в пространстве состояний содержит
интегратор ($\int$) для перехода от $\dot{\vect{x}}$ к $\vect{x}$,
блоки умножения на матрицы $\mat{A}$, $\mat{B}$, $\mat{C}$, $\mat{D}$
и сумматоры.
\end{notebox}

\subsection{Применение в проекте SAMURAI}

Каждый фильтр Калмана в проекте является реализацией модели
в пространстве состояний:

\begin{center}
\begin{tabular}{lll}
\toprule
\textbf{Модуль} & \textbf{Вектор состояния $\vect{x}$} & \textbf{Размерность $n$}\\
\midrule
Velocity EKF (гл.~\ref{ch:velocity_ekf}) & $(v_x, v_y)\T$ & 2\\
IMU EKF (гл.~\ref{ch:imu_ekf}) & $(\psi, b_z)\T$ & 2\\
AccelPosition (гл.~\ref{ch:accel_position}) & $(x, y, v_x, v_y)\T$ & 4\\
Одометрия (гл.~\ref{ch:odometry}) & $(x, y, \theta)\T$ & 3\\
\bottomrule
\end{tabular}
\end{center}

% ═══════════════════════════════════════════════════════════════
\section{Линеаризация нелинейных систем}
\label{sec:linearization}

Реальные объекты, включая наземного робота, описываются
\textbf{нелинейными} уравнениями:

\begin{equation}
  \dot{\vect{x}}(t) = \vect{f}(\vect{x}, \vect{u}), \qquad
  \vect{y}(t) = \vect{g}(\vect{x}, \vect{u})
  \label{eq:nonlinear_system}
\end{equation}

\subsection{Пример: транспортная тележка}

Модель движения мобильного робота (транспортной тележки) описывается
системой шести дифференциальных уравнений:

\begin{align}
  \dot{\theta} &= -\frac{1}{T_1}\theta + \frac{1}{T_1}u_1 \label{eq:cart_theta}\\
  \dot{f} &= -\frac{1}{T_2}f + \frac{1}{T_2}u_2 \label{eq:cart_force}\\
  \dot{v} &= -\frac{\alpha}{m}v + \frac{1}{m}f \label{eq:cart_vel}\\
  \dot{x} &= v\cos(\psi + \theta) \label{eq:cart_x}\\
  \dot{\psi} &= \frac{1}{L}v\sin(\theta) \label{eq:cart_psi}\\
  \dot{z} &= -v\sin(\psi + \theta) \label{eq:cart_z}
\end{align}

где $\theta$ --- угол поворота колёс, $\psi$ --- угол рыскания,
$f$ --- движущая сила, $v$ --- скорость, $x,z$ --- координаты;
$T_1, T_2, \alpha, m, L$ --- параметры объекта.

Уравнения~\eqref{eq:cart_x}--\eqref{eq:cart_z} \textbf{нелинейны}
(содержат $\sin$ и $\cos$), что делает систему нелинейной.

\begin{notebox}
Модель робота SAMURAI использует ту же структуру:
уравнения одометрии (гл.~\ref{ch:odometry}) $\dot{x} = v\cos\theta$,
$\dot{y} = v\sin\theta$ нелинейны из-за тригонометрических функций
от угла курса $\theta$.
\end{notebox}

\subsection{Метод линеаризации по Тейлору (матрица Якоби)}
\label{sec:jacobian}

Линеаризация проводится в окрестности \textbf{опорного режима}
$(\bar{\vect{x}}, \bar{\vect{u}})$, в котором $\vect{f}(\bar{\vect{x}}, \bar{\vect{u}}) = \vect{0}$.

Разложение в ряд Тейлора с отбрасыванием членов высших порядков:

\begin{equation}
  \boxed{
  \Delta\dot{\vect{x}} = \mat{A}\,\Delta\vect{x} + \mat{B}\,\Delta\vect{u}, \qquad
  \Delta\vect{y} = \mat{C}\,\Delta\vect{x} + \mat{D}\,\Delta\vect{u}
  }
  \label{eq:linearized}
\end{equation}

где матрицы --- \textbf{якобианы}, вычисленные в опорной точке:

\begin{equation}
  \mat{A} = \left.\frac{\partial \vect{f}}{\partial \vect{x}}\right|_{\bar{\vect{x}},\bar{\vect{u}}}, \quad
  \mat{B} = \left.\frac{\partial \vect{f}}{\partial \vect{u}}\right|_{\bar{\vect{x}},\bar{\vect{u}}}, \quad
  \mat{C} = \left.\frac{\partial \vect{g}}{\partial \vect{x}}\right|_{\bar{\vect{x}},\bar{\vect{u}}}, \quad
  \mat{D} = \left.\frac{\partial \vect{g}}{\partial \vect{u}}\right|_{\bar{\vect{x}},\bar{\vect{u}}}
  \label{eq:jacobians}
\end{equation}

\subsection{Пример: якобиан для транспортной тележки}

Для опорного режима: $v(t) = v^0$, $\theta(t) = 0$, $\psi(t) = 0$:

\begin{equation}
  \mat{A} = \begin{pmatrix}
    -1/T_1 & 0 & 0 & 0 & 0 & 0\\
    0 & -1/T_2 & 0 & 0 & 0 & 0\\
    0 & 1/m & -\alpha/m & 0 & 0 & 0\\
    0 & 0 & 1 & 0 & 0 & 0\\
    v^0/L & 0 & 0 & 0 & 0 & 0\\
    -v^0 & 0 & 0 & 0 & -v^0 & 0
  \end{pmatrix}
  \label{eq:cart_A}
\end{equation}

\subsection{Связь с EKF в проекте SAMURAI}

\textbf{Расширенный фильтр Калмана (EKF)} --- это именно
\textbf{линеаризация на каждом шаге}. Матрица перехода $\mat{F}$
в фильтре Калмана (гл.~\ref{ch:imu_ekf}, формула~\eqref{eq:F_matrix})
--- это якобиан нелинейной модели:

\begin{equation}
  \mat{F}_k = \left.\frac{\partial \vect{f}}{\partial \vect{x}}\right|_{\vect{x}_{k-1}}
  = \begin{pmatrix} 1 & -\Delta t \\ 0 & 1 \end{pmatrix}
  \label{eq:ekf_jacobian}
\end{equation}

Здесь нелинейная модель $\psi_k = \psi_{k-1} + (g_z - b_z)\Delta t$
линеаризуется относительно текущего состояния $(\psi_{k-1}, b_{z,k-1})$,
а якобиан $\mat{F}$ содержит производные $\partial\dot{\psi}/\partial\psi = 1$,
$\partial\dot{\psi}/\partial b_z = -\Delta t$.

% ═══════════════════════════════════════════════════════════════
\section{Передаточная функция}
\label{sec:transfer_func}

Для линейной стационарной системы~\eqref{eq:state_space}
\textbf{передаточная функция} определяется через преобразование Лапласа:

\begin{equation}
  \boxed{W(p) = \mat{C}(p\mat{E}_n - \mat{A})^{-1}\mat{B} + \mat{D}}
  \label{eq:transfer_func}
\end{equation}

где $p = \frac{d}{dt}$ --- оператор дифференцирования, $\mat{E}_n$ --- единичная матрица.

\subsection{Характеристический полином и резольвента}

Резольвента матрицы $\mat{A}$:

\begin{equation}
  (p\mat{E}_n - \mat{A})^{-1} = \frac{1}{\chi_n(p)}\widetilde{\mat{B}}(p)
  \label{eq:resolvent}
\end{equation}

где $\chi_n(p) = \det(p\mat{E}_n - \mat{A}) = p^n + a_1 p^{n-1} + \ldots + a_n$
--- характеристический полином, а $\widetilde{\mat{B}}(p)$ --- присоединённая матрица.

\subsection{Алгоритм Фаддеева--Леверье}
\label{sec:faddeev}

Коэффициенты характеристического полинома и матрицы
$\widetilde{\mat{B}}_k$ вычисляются рекуррентно:

\begin{equation}
  \boxed{
  \begin{aligned}
    \widetilde{\mat{B}}_1 &= \mat{E}_n, \quad & a_1 &= -\mathrm{tr}(\mat{A}\widetilde{\mat{B}}_1)\\
    \widetilde{\mat{B}}_k &= \mat{A}\widetilde{\mat{B}}_{k-1} + a_{k-1}\mat{E}_n, \quad
    & a_k &= -\frac{1}{k}\mathrm{tr}(\mat{A}\widetilde{\mat{B}}_k)\\
    & & \multicolumn{2}{l}{k = 2, \ldots, n}
  \end{aligned}
  }
  \label{eq:faddeev}
\end{equation}

Проверка: $\mat{A}\widetilde{\mat{B}}_n + a_n\mat{E}_n = \mat{0}$
(теорема Гамильтона--Кели).

\subsection{Уравнение <<вход--выход>>}

Передаточная функция связывает выход $y$ с входом $u$ через
дифференциальное уравнение:

\begin{equation}
  p^n y + a_1 p^{n-1} y + \ldots + a_n y = b_0 p^n u + b_1 p^{n-1} u + \ldots + b_n u
  \label{eq:io_equation}
\end{equation}

\begin{notebox}
В проекте SAMURAI фильтры Калмана работают в \textbf{дискретном}
пространстве состояний:
$\vect{x}_{k+1} = \mat{F}\,\vect{x}_k + \mat{B}\,\vect{u}_k$,
что является дискретным аналогом непрерывной модели~\eqref{eq:state_space}.
Переход к дискретной модели:
$\mat{A}_d = e^{\mat{A}T}$, $\mat{B}_d = \int_0^T e^{\mat{A}\tau}d\tau\,\mat{B}$.
\end{notebox}

% ═══════════════════════════════════════════════════════════════
\section{Устойчивость систем управления}
\label{sec:stability}

\subsection{Определение устойчивости}

Система~\eqref{eq:state_space} \textbf{асимптотически устойчива},
если решение $\vect{x}(t) \to \vect{0}$ при $t \to \infty$
для любых начальных условий $\vect{x}(0)$.

\subsection{Критерий собственных значений}

Система устойчива тогда и только тогда, когда все собственные числа
матрицы $\mat{A}$ имеют отрицательные вещественные части:

\begin{equation}
  \boxed{\mathrm{Re}(\lambda_i) < 0, \quad \forall\, i = 1, \ldots, n}
  \label{eq:stability_eigenvalues}
\end{equation}

где $\lambda_i$ --- корни характеристического уравнения $\det(\lambda\mat{E} - \mat{A}) = 0$.

\subsection{Критерий Гурвица}

Для характеристического полинома $\chi(p) = p^n + a_1 p^{n-1} + \ldots + a_n$
формируется матрица Гурвица:

\begin{equation}
  \mat{H} = \begin{pmatrix}
    a_1 & a_3 & a_5 & \cdots\\
    1   & a_2 & a_4 & \cdots\\
    0   & a_1 & a_3 & \cdots\\
    0   & 1   & a_2 & \cdots\\
    \vdots & & & \ddots
  \end{pmatrix}
  \label{eq:hurwitz}
\end{equation}

\textbf{Условие устойчивости}: все угловые миноры матрицы $\mat{H}$
положительны:

\begin{equation}
  \Delta_1 = a_1 > 0, \quad
  \Delta_2 = a_1 a_2 - a_3 > 0, \quad
  \ldots, \quad \Delta_n > 0
  \label{eq:hurwitz_cond}
\end{equation}

\subsection{Решение уравнения состояния}

Аналитическое решение линейной системы (формула Коши):

\begin{equation}
  \vect{x}(t) = e^{\mat{A}(t - t_0)}\vect{x}(t_0) + \int_{t_0}^{t} e^{\mat{A}(t - \tau)}\mat{B}\,\vect{u}(\tau)\,d\tau
  \label{eq:cauchy}
\end{equation}

где $e^{\mat{A}t}$ --- матричная экспонента. Первое слагаемое ---
свободное движение (от начальных условий), второе ---
вынужденное движение (от входного воздействия).

\subsection{Жорданова форма и матричная экспонента}

Через замену $\vect{x} = \mat{S}\vect{z}$ система приводится к жордановой форме:

\begin{equation}
  \dot{\vect{z}} = \mat{J}\vect{z} + \mat{S}^{-1}\mat{B}\vect{u}, \quad
  \mat{J} = \mat{S}^{-1}\mat{A}\mat{S}
  \label{eq:jordan}
\end{equation}

Для диагональной жордановой формы экспонента:

\begin{equation}
  e^{\mat{J}t} = \mathrm{diag}\!\left(e^{\lambda_1 t}, e^{\lambda_2 t}, \ldots, e^{\lambda_n t}\right)
  \label{eq:matrix_exp}
\end{equation}

Отсюда видно: при $\mathrm{Re}(\lambda_i) < 0$ каждая мода затухает
экспоненциально, что обеспечивает устойчивость.

% ═══════════════════════════════════════════════════════════════
\section{Дискретные системы и фильтр Калмана}
\label{sec:discrete_kalman}

\subsection{Дискретная модель в пространстве состояний}

\begin{equation}
  \vect{x}_{k+1} = \mat{A}_d\,\vect{x}_k + \mat{B}_d\,\vect{u}_k, \qquad
  \vect{y}_k = \mat{C}_d\,\vect{x}_k + \mat{D}_d\,\vect{u}_k
  \label{eq:discrete_ss}
\end{equation}

где $\mat{A}_d = e^{\mat{A}T}$, $T$ --- период дискретизации.

\subsection{Общие уравнения фильтра Калмана}

Фильтр Калмана решает задачу оптимальной оценки состояния
стохастической системы:

\begin{equation}
  \vect{x}_{k} = \mat{F}\,\vect{x}_{k-1} + \mat{B}\,\vect{u}_{k-1} + \vect{w}_{k-1}, \qquad
  \vect{z}_k = \mat{H}\,\vect{x}_k + \vect{v}_k
  \label{eq:kalman_model}
\end{equation}

где $\vect{w} \sim \mathcal{N}(\vect{0}, \mat{Q})$ --- процессный шум,
$\vect{v} \sim \mathcal{N}(\vect{0}, \mat{R})$ --- шум измерений.

\textbf{Шаг предсказания (Predict):}

\begin{align}
  \hat{\vect{x}}_k^- &= \mat{F}\,\hat{\vect{x}}_{k-1} + \mat{B}\,\vect{u}_{k-1}
  \label{eq:kf_predict_x}\\
  \mat{P}_k^- &= \mat{F}\,\mat{P}_{k-1}\,\mat{F}\T + \mat{Q}
  \label{eq:kf_predict_P}
\end{align}

\textbf{Шаг коррекции (Update):}

\begin{align}
  \vect{y}_k &= \vect{z}_k - \mat{H}\,\hat{\vect{x}}_k^- && \text{(инновация)}
  \label{eq:kf_innovation}\\
  \mat{S}_k &= \mat{H}\,\mat{P}_k^-\,\mat{H}\T + \mat{R} && \text{(ковариация инновации)}
  \label{eq:kf_S}\\
  \mat{K}_k &= \mat{P}_k^-\,\mat{H}\T\,\mat{S}_k^{-1} && \text{(усиление Калмана)}
  \label{eq:kf_gain}\\
  \hat{\vect{x}}_k &= \hat{\vect{x}}_k^- + \mat{K}_k\,\vect{y}_k && \text{(обновление состояния)}
  \label{eq:kf_update_x}\\
  \mat{P}_k &= (\mat{I} - \mat{K}_k\mat{H})\,\mat{P}_k^- && \text{(обновление ковариации)}
  \label{eq:kf_update_P}
\end{align}

\begin{formulabox}
Фильтр Калмана \textbf{минимизирует} дисперсию ошибки оценки
$E[\|\vect{x} - \hat{\vect{x}}\|^2]$. Усиление $\mat{K}$ автоматически
взвешивает доверие между предсказанием (модель) и измерением (датчик):
\begin{itemize}
  \item $\mat{R} \ll \mat{P}^-$: $\mat{K} \to \mat{H}^{-1}$ --- доверяем датчику
  \item $\mat{R} \gg \mat{P}^-$: $\mat{K} \to \mat{0}$ --- доверяем модели
\end{itemize}
\end{formulabox}

\subsection{Расширенный фильтр Калмана (EKF)}

Для нелинейной системы $\vect{x}_k = \vect{f}(\vect{x}_{k-1}, \vect{u}_{k-1})$,
$\vect{z}_k = \vect{h}(\vect{x}_k)$ используется \textbf{линеаризация
на каждом шаге}:

\begin{equation}
  \mat{F}_k = \left.\frac{\partial \vect{f}}{\partial \vect{x}}\right|_{\hat{\vect{x}}_{k-1}}, \qquad
  \mat{H}_k = \left.\frac{\partial \vect{h}}{\partial \vect{x}}\right|_{\hat{\vect{x}}_k^-}
  \label{eq:ekf_linearize}
\end{equation}

Затем якобианы $\mat{F}_k$, $\mat{H}_k$ подставляются в стандартные
уравнения фильтра Калмана~\eqref{eq:kf_predict_P}--\eqref{eq:kf_update_P}.

% ═══════════════════════════════════════════════════════════════
\section{Модель электродвигателя постоянного тока}
\label{sec:dc_motor}

Двигатели робота SAMURAI управляются ШИМ-сигналами через драйвер TB6612.
Электромеханические процессы в двигателе постоянного тока описываются
дифференциальным уравнением второго порядка:

\begin{equation}
  \boxed{T_a^2\ddot{\omega} + T_y\dot{\omega} + \omega = k\,u}
  \label{eq:dc_motor}
\end{equation}

где:
\begin{itemize}
  \item $\omega$ --- угловая скорость вала двигателя
  \item $u$ --- управляющее воздействие (напряжение / ШИМ)
  \item $T_a$ --- электромеханическая постоянная времени
  \item $T_y$ --- электрическая постоянная времени
  \item $k$ --- коэффициент передачи
\end{itemize}

\subsection{Пространство состояний для двигателя}

Вводя $x_1 = \omega$, $x_2 = \dot{\omega}$:

\begin{equation}
  \begin{pmatrix} \dot{x}_1 \\ \dot{x}_2 \end{pmatrix}
  = \begin{pmatrix} 0 & 1 \\ -1/T_a^2 & -T_y/T_a^2 \end{pmatrix}
  \begin{pmatrix} x_1 \\ x_2 \end{pmatrix}
  + \begin{pmatrix} 0 \\ k/T_a^2 \end{pmatrix} u
  \label{eq:motor_ss}
\end{equation}

Характеристический полином:

\begin{equation}
  \chi(\lambda) = \lambda^2 + \frac{T_y}{T_a^2}\lambda + \frac{1}{T_a^2}
  \label{eq:motor_char}
\end{equation}

\textbf{Условие устойчивости} (по Гурвицу): $T_y > 0$, $T_a > 0$ ---
выполняется для всех реальных двигателей.

\subsection{Передаточная функция двигателя}

\begin{equation}
  W(p) = \frac{k}{T_a^2 p^2 + T_y p + 1}
  \label{eq:motor_tf}
\end{equation}

При $T_a \ll T_y$ (типично для малых моторов) двигатель приближённо
описывается звеном первого порядка:

\begin{equation}
  W(p) \approx \frac{k}{T_y p + 1}
  \label{eq:motor_tf_approx}
\end{equation}

\begin{notebox}
В проекте SAMURAI ШИМ-управление (гл.~\ref{ch:pwm_motor})
работает с нормированным throttle $t \in [-1, 1]$, что соответствует
$u/u_{\max}$. Драйвер TB6612 в режиме Slow Decay обеспечивает
линейную зависимость $\omega \propto t$ в установившемся режиме.
\end{notebox}

% ═══════════════════════════════════════════════════════════════
\section{Пропорциональный регулятор (P-регулятор)}
\label{sec:p_controller}

\subsection{Закон управления}

\begin{equation}
  u(t) = K_P \cdot e(t), \qquad e(t) = r(t) - y(t)
  \label{eq:p_controller}
\end{equation}

где $r(t)$ --- задание, $y(t)$ --- выход, $K_P$ --- коэффициент усиления.

\subsection{Передаточная функция замкнутой системы}

Для объекта с передаточной функцией $W_{\text{obj}}(p)$
и P-регулятора с коэффициентом $K_P$:

\begin{equation}
  W_{\text{зс}}(p) = \frac{K_P \cdot W_{\text{obj}}(p)}{1 + K_P \cdot W_{\text{obj}}(p)}
  \label{eq:closed_loop}
\end{equation}

\subsection{Статическая ошибка}

P-регулятор не может полностью устранить установившуюся ошибку
при ступенчатом воздействии:

\begin{equation}
  e_{\text{уст}} = \frac{1}{1 + K_P \cdot k_{\text{obj}}}
  \label{eq:static_error}
\end{equation}

Увеличение $K_P$ снижает ошибку, но может привести к
потере устойчивости.

\subsection{Применение P-регуляторов в проекте}

\begin{center}
\begin{tabular}{llll}
\toprule
\textbf{Модуль} & \textbf{Управляемая переменная} & $K_P$ & \textbf{Ограничение}\\
\midrule
Pure Pursuit (гл.~\ref{ch:pure_pursuit}) & Угловая скорость $\omega$ & $2.5$ & $|\omega| \leq 0.8$\\
Follow-Me (гл.~\ref{ch:follow_me}) & Угловая скорость $\omega$ & $0.8$ & ---\\
Follow-Me (гл.~\ref{ch:follow_me}) & Линейная скорость $v$ & $0.3$ & $v \leq 0.20$\\
Patrol (гл.~\ref{ch:patrol}) & Целевая позиция & --- & $d < 0.20\,\text{м}$\\
\bottomrule
\end{tabular}
\end{center}

% ═══════════════════════════════════════════════════════════════
\section{Переходная характеристика}
\label{sec:transient}

Переходная характеристика $h(t)$ --- реакция системы на единичное
ступенчатое воздействие $u(t) = 1(t)$ при нулевых начальных условиях:

\begin{equation}
  h(t) = \mat{C}\int_0^t e^{\mat{A}\tau}\,d\tau\,\mat{B} + \mat{D}
  \label{eq:step_response}
\end{equation}

При отсутствии нулевых собственных значений у $\mat{A}$:

\begin{equation}
  h(t) = \mat{C}\left(e^{\mat{A}t} - \mat{E}\right)\mat{A}^{-1}\mat{B} + \mat{D}
  \label{eq:step_response2}
\end{equation}

Для системы второго порядка с $\omega_n$ (собственная частота)
и $\zeta$ (коэффициент демпфирования):

\begin{equation}
  h(t) = 1 - \frac{e^{-\zeta\omega_n t}}{\sqrt{1-\zeta^2}}
  \sin\!\left(\omega_n\sqrt{1-\zeta^2}\,t + \arccos\zeta\right)
  \label{eq:2nd_order_response}
\end{equation}

\begin{notebox}
Косинусное масштабирование скорости в Pure Pursuit
(гл.~\ref{ch:pure_pursuit}, формула~\eqref{eq:cos_factor})
обеспечивает плавную переходную характеристику без перерегулирования
при развороте робота.
\end{notebox}

% ═══════════════════════════════════════════════════════════════
\section{Управляемость и наблюдаемость}
\label{sec:controllability}

\subsection{Матрица управляемости}

\begin{equation}
  \mat{S}_v = \bigl[\mat{B},\; \mat{A}\mat{B},\; \mat{A}^2\mat{B},\; \ldots,\; \mat{A}^{n-1}\mat{B}\bigr]
  \label{eq:controllability}
\end{equation}

Система полностью управляема $\Leftrightarrow$ $\mathrm{rank}(\mat{S}_v) = n$.

\subsection{Матрица наблюдаемости}

\begin{equation}
  \mat{S}_o = \begin{pmatrix} \mat{C} \\ \mat{C}\mat{A} \\ \vdots \\ \mat{C}\mat{A}^{n-1} \end{pmatrix}
  \label{eq:observability}
\end{equation}

Система полностью наблюдаема $\Leftrightarrow$ $\mathrm{rank}(\mat{S}_o) = n$.

\begin{notebox}
Фильтр Калмана может оценивать состояние только
\textbf{наблюдаемой} системы. EKF для угла курса
(гл.~\ref{ch:imu_ekf}) наблюдаем, так как гироскоп
измеряет $\omega_z = \dot{\psi} - b_z$, что позволяет
оценить и $\psi$, и $b_z$.
\end{notebox}

% ═══════════════════════════════════════════════════════════════
\section{Сводная таблица: теория $\to$ реализация}

\begin{center}
\begin{tabular}{lll}
\toprule
\textbf{Теоретический концепт} & \textbf{Глава} & \textbf{Реализация}\\
\midrule
$\dot{\vect{x}} = \mat{A}\vect{x} + \mat{B}\vect{u}$ & \ref{ch:velocity_ekf}, \ref{ch:imu_ekf}
  & Predict-шаг EKF\\
Якобиан $\mat{F} = \partial\vect{f}/\partial\vect{x}$ & \ref{ch:imu_ekf}
  & Матрица перехода $\mat{F}$\\
$\mat{K} = \mat{P}\mat{H}\T\mat{S}^{-1}$ & \ref{ch:velocity_ekf}
  & Усиление Калмана\\
$u = K_P \cdot e$ & \ref{ch:pure_pursuit}, \ref{ch:follow_me}
  & Угловой P-регулятор\\
$T_a^2\ddot{\omega} + T_y\dot{\omega} = ku$ & \ref{ch:pwm_motor}
  & ШИМ-управление мотором\\
$\mathrm{Re}(\lambda_i) < 0$ & \ref{ch:imu_ekf}
  & Устойчивость EKF\\
$e^{\mat{A}t}$ & \ref{ch:accel_position}
  & Интегрирование позиции\\
$\vect{x}_{k+1} = \mat{A}_d\vect{x}_k$ & \ref{ch:odometry}
  & Дискретная одометрия\\
\bottomrule
\end{tabular}
\end{center}
